\section*{第 \ref{ch:b2:flow-balances} 章——流动中的动量与能量平衡}

\begin{solution}{exo:b2:flow-balances:1}
$s = \qty{7.1}{cm^2}$, $v = 0.025/7.1 \times 10^{-4} = \qty{35}{m/s}$; $F = D_mv = 25
\times 35 = \qty{880}{N}$ 作用于喷嘴和墙面上；若墙面后退：$\rho s(v -
u)^2 = 1000 \times 7.1 \times 10^{-4} \times 25.4^2 = \qty{460}{N}$.
\end{solution}

\begin{solution}{exo:b2:flow-balances:2}
$F = 250 \times 3000 = \qty{750}{kN}$；$a_0 = 750/60 - 9.81 = \qty{2.7}{m/s^2}$；在
\qty{20}{t}时：$37.5 - 9.8 = \qty{28}{m/s^2}$。$\Delta v = 3000\ln3 = \qty{3.3}{km/s}$ 在
自由空间中；减去 $g\tau = \qty{1.6}{km/s}$：\qty{1.7}{km/s}。
\end{solution}

\begin{solution}{exo:b2:flow-balances:3}
$S = \qty{177}{cm^2}$：$PS = \qty{7.1}{kN}$，$\rho Sv^2 = \qty{71}{N}$。弯管所受力的每个分量
为 $PS + \rho Sv^2 = \qty{7.1}{kN}$：\qty{10}{kN}，沿
外角平分线方向，其中 $99\%$ 来自压强。开口出口（表压
为零）：$7.1$ kN 沿入口方向，\qty{71}{N} 沿
出口方向：\qty{7.1}{kN}，几乎沿入口轴线方向。
\end{solution}

\begin{solution}{exo:b2:flow-balances:4}
$F = 80 \times 550 = \qty{44}{kN}$。飞行时：$F = 80(550 - 240) = \qty{25}{kN}$，$Fv =
\qty{6.0}{MW}$；$\mathcal P_{\text{动能}} = \tfrac12 \times 80 \times (550^2 - 240^2) =
\qty{9.8}{MW}$; $\eta_p = 0.61 = 2 \times 240/790$.
\end{solution}

\begin{solution}{exo:b2:flow-balances:5}
$D_m = \qty{500}{kg/s}$. (a) $F = 1000(60 - u)$ N, $\mathcal P = 1000u(60 - u)$; $u =
\qty{30}{m/s}$, $\omega = 37.5$ rad/s $= \qty{360}{rpm}$. (b) \qty{900}{kW} $=
\tfrac12D_mv^2$: $100\%$. (c) $F = D_m(v - u)(1 + \cos15^\circ) = \qty{29.5}{kN}$, $\mathcal P
= \qty{885}{kW}$，$98\%$。(d) $\Gamma = FR = \qty{24}{kN.m}$；静止时力
加倍（$\Gamma = \qty{48}{kN.m}$）；当 $u = v$ 时，力与力矩均为零。
\end{solution}

\begin{solution}{exo:b2:flow-balances:6}
(a) 每个喷嘴 $s = \qty{7.1}{mm^2}$，流量各为 \qty{0.10}{L/s}：$w = \qty{14}{m/s}$。(b)
$\Gamma_0 = D_mRw = 0.2 \times 0.15 \times 14 = \qty{0.42}{N.m}$. (c) $D_mR(w - R\omega) =
5 \times 10^{-3}$: $w - R\omega = \qty{0.17}{m/s}$, $\omega = \qty{93}{rad/s} \approx \qty{890}{rpm}$;
无摩擦时 $R\omega = w$：\qty{900}{rpm}。(d) \qty{0.17}{m/s} 反向，随后
为零：水竖直下落。
\end{solution}

\begin{solution}{exo:b2:flow-balances:7}
(a) $D_m = \rho Av_i$，动量通量 $D_m\cdot2v_i$：$T = 2\rho Av_i^2$；$A = \qty{78.5}{m^2}$，
$T = \qty{19.6}{kN}$: $v_i = \qty{10}{m/s}$. (b) $\mathcal P = \tfrac12D_m(2v_i)^2 = Tv_i =
\qty{200}{kW}$。(c) $\mathcal P = T^{3/2}/\sqrt{2\rho A}$：面积加倍时
使诱导功率除以 $\sqrt2$——较大的转子推动更多
空气，但速度更慢。(d) $A = \qty{0.20}{m^2}$，$v_i = \sqrt{9.81/0.47} = \qty{4.6}{m/s}$，$\mathcal P
= \qty{45}{W}$：\qty{45}{W/kg}，与 \qty{100}{W/kg} 相比——盘面载荷 $T/A$
较低。
\end{solution}

\begin{solution}{exo:b2:flow-balances:8}
(a) $v_2 = \qty{1.0}{m/s}$；伯努利给出 $\tfrac12\rho(16 - 1) = \qty{7.5}{kPa}$；实际
$\rho v_2(v_1 - v_2) = \qty{3.0}{kPa}$；损失 \qty{4.5}{kPa} $= \qty{0.46}{m}$。(b) 回收 $6$ kPa，
损失 \qty{1.5}{kPa}。(c) $\tfrac12\rho v_1^2 = \qty{8}{kPa}$，\qty{0.82}{m}。(d)
$D_V = \qty{7.9}{L/s}$：\qty{35}{W} 和 \qty{63}{W}。
\end{solution}

\begin{solution}{exo:b2:flow-balances:9}
(a) $v = 0.01/2.83 \times 10^{-3} = \qty{3.5}{m/s}$，$v^2/2g = \qty{0.64}{m}$；摩擦损失
$0.025 \times 1000 \times 0.64 = \qty{16}{m}$，管件 \qty{1.3}{m}：\qty{17.3}{m}。(b)
$H = 30 + 17.3 + 0.64 = \qty{48}{m}$；$\rho gD_VH = \qty{4.7}{kW}$；轴处功率为
\qty{7.2}{kW}。(c) 进口（不计吸入损失）：$P_0 - \rho g(5 + 0.64) = \qty{45}{kPa}$；
出口 $45 + 470 = \qty{515}{kPa}$。(d) $P_{\text{进口}} \ge \qty{2.3}{kPa}$：$h \le (100 -
2.3)/9.81 - 0.64 = \qty{9.3}{m}$（计入吸入损失时更小；实际中为 $7$--\qty{8}{m}
）。
\end{solution}

\begin{solution}{exo:b2:flow-balances:10}
(a) $v = \sqrt{2 \times 4 \times 10^5/1000} = \qty{28}{m/s}$; $s = \qty{3.8}{cm^2}$, $D_m =
\qty{10.8}{kg/s}$, $F = D_mv = 2\Delta P\,s = \qty{300}{N}$; $a = 300/1.1 - 9.8 =
\qty{260}{m/s^2}$。(b) 空气体积为 \qty{1.5}{L} 时：$P = \qty{3.3}{bar}$，$v = \sqrt{2 \times 2.3 \times
10^5/1000} = \qty{22}{m/s}$。(c) 平均出口速率 $\approx \qty{22}{m/s}$：燃烧时长 $\approx 10^{-3}
/(3.8 \times 10^{-4} \times 22) = \qty{0.12}{s}$；用火箭方程，取 $v_e \approx \qty{22}{m/s}$：
$\Delta v \approx 22\ln(1.1/0.1) = \qty{50}{m/s}$。(d) 剩余的 \qty{2.5}{bar} 空气
排出，并略微增加推力；水过多则留下很少的空气，
压强过早下降；水过少则可抛出的质量很少——
约占体积的三分之一最佳。
\end{solution}

\begin{solution}{exo:b2:flow-balances:11}
(a) $\Gamma = D_mr_1v_{\theta1} = 3 \times 10^4 \times 20 = \qty{600}{kN.m}$, $\omega = 15.7$ rad/s,
$\mathcal P = \qty{9.4}{MW}$. (b) $H = \mathcal P/\rho gD_V = \qty{32}{m}$. (c) $\Gamma = 3 \times
10^4(20 - 1.5) = \qty{555}{kN.m}$，$\mathcal P = \qty{11.6}{MW}$（水头 \qty{39}{m}）；
出口涡流携带 $v_{\theta2}^2/2g = \qty{0.46}{m}$，约 $1\%$。(d) 佩尔顿射流
处于大气压强：需较高水头才能获得高速，其
流量受喷管限制；弗朗西斯转轮充满有压的水，
并可在中等水头下通过大流量。
\end{solution}

\begin{solution}{exo:b2:flow-balances:12}
(a) $v_1h_1 = v_2h_2 = q$。(b) 作用于截面上的单位宽度静水压力
$\int_0^h\rho g(h - z)\dd z = \tfrac12\rho gh^2$；动量：$\tfrac12\rho g(h_1^2
- h_2^2) = \rho q(v_2 - v_1)$；除以 $\rho$，并利用 $q = v_1h_1 = v_2h_2$。
(c) 取 $v_2 = v_1h_1/h_2$，且 $h_1 \ne h_2$：$\tfrac12g(h_1 + h_2)h_2 = v_1^2h_1$，
即 $(h_2/h_1)^2 + h_2/h_1 - 2\mathrm{Fr}_1^2 = 0$；$h_2 > h_1$ 需要 $\mathrm{Fr}_1 > 1$。
(d) $\Delta H = (h_1 + v_1^2/2g) - (h_2 + v_2^2/2g) = (h_2 - h_1)^3/4h_1h_2$.
$\mathrm{Fr}_1 = 2.86$, $h_2/h_1 = 3.57$, $h_2 = \qty{0.71}{m}$, $v_2 = \qty{1.1}{m/s}$,
$\Delta H = 0.514^3/0.57 = \qty{0.24}{m}$；$\rho gq\,\Delta H = \qty{1.9}{kW}$ 每米
宽度。
\end{solution}

\begin{solution}{pb:b2:flow-balances:1}
\textbf{1.} $F = (D_a + D_f)v_e = 51 \times 600 = \qty{30.6}{kN}$.

\textbf{2.} $F = 51 \times 600 - 50 \times 250 = \qty{18.1}{kN}$; $Fv = \qty{4.5}{MW}$.

\textbf{3.} $\mathcal P_{\text{动能}} = \tfrac12 \times 51 \times 600^2 - \tfrac12 \times 50 \times 250^2 =
\qty{7.6}{MW}$；$\eta_p = 0.59$。取 $D_f = 0$ 时：$\eta_p = 2v(v_e - v)/(v_e^2 - v^2)
= 2v/(v_e + v) = 500/850$。

\textbf{4.} 热量输入 \qty{43}{MW}：热效率 $7.6/43 = 18\%$；总效率 $4.5/43 = 10\%$。

\textbf{5.} $\tfrac12 \times 500(v_e'^2 - 250^2) = 7.6 \times 10^6$，得 $v_e' = \qty{305}{m/s}$；
$F = 500 \times 55 = \qty{27.5}{kN}$；$\eta_p = 500/555 = 0.90$。推动更多空气、使其
速度更低，在相同功率下可获得更大推力，因此采用风扇。

\textbf{6.} 它既需要外界空气作为工质，又需要其作为氧化剂；
火箭则自身携带二者。

\textbf{7.} 空气以 \qty{70}{m/s} 的速度进入（在发动机参考系中向后），
并以向前的分量 $150\cos45^\circ = \qty{106}{m/s}$ 离开，
空气获得的动量向前，因而制动力为 $500(106 + 70) = \qty{88}{kN}$；
若同样的气流向后喷出，则推力为 $500(150 - 70) =
\qty{40}{kN}$。

\textbf{8.} 在 $\dd t$ 时间内，系统（火箭质量 $m$，加上质量 $D_m\dd t$ 的气体，气体以 $v - v_e$ 喷出）满足
$(m - D_m\dd t)(v + \dd v) + D_m\dd t(v - v_e) - mv = -mg\,\dd t$，
由此 $m\,\dd v/\dd t = D_mv_e - mg$。

\textbf{9.} $D_m = \qty{1333}{kg/s}$；$F = \qty{4.0}{MN}$；$a = 13.3 - 9.8 =
\qty{3.5}{m/s^2}$ 起飞时，$40 - 9.8 = \qty{30}{m/s^2}$ 燃料耗尽时。

\textbf{10.} $v = v_e\ln(m_0/m) - gt$：$3000\ln3 - 9.81 \times 150 = 3296 - 1472 =
\qty{1.8}{km/s}$；重力损失为 \qty{1.5}{km/s}。

\textbf{11.} 第一级：$3000\ln(300/200) = \qty{1.2}{km/s}$；抛去 \qty{20}{t}；第二级
：$3000\ln(180/80) = \qty{2.4}{km/s}$，合计 \qty{3.6}{km/s}，而单级火箭仅为 \qty{3.3}{km/s}
（计入重力后，差距更大）。

\textbf{12.} $v_e = 7800/\ln3 = \qty{7.1}{km/s}$：没有任何化学推进剂能达到此值
（氢–氧推进剂仅给出 \qty{4.5}{km/s}）；必须采用多级火箭。

\textbf{13.} 喷流的竖直动量通量 $D_mv = \qty{4.0}{MN}$ 被去除，
并产生一个大小相等的水平动量通量：即竖直向下的 \qty{4}{MN} 力，
加上水平方向的 \qty{4}{MN}，合计 \qty{5.7}{MN}。

\textbf{14.} 发动机的动量平衡在出口截面上包含 $(P_e - P_0)S_e$；
由于 $P_0$ 随海拔升高而下降，该项随之增大，因此真空中推力
约增加 $10$--$15\%$。

\textbf{15.} 每个水轮的 $D_m = \qty{3e4}{kg/s}$：$F = D_m(v - u)(1 + \cos15^\circ) =
1.97 \times 3 \times 10^4(62.6 - u)$；$\mathcal P = Fu$，在 $u = v/2 = \qty{31.3}{m/s}$ 时最大。

\textbf{16.} $\omega = u/R = 20.9$ rad/s $= \qty{199}{rpm}$；$F = \qty{1.85}{MN}$，$\Gamma = FR
= \qty{2.8}{MN.m}$。

\textbf{17.} $f = pn/60$：$p = 3000/199 \approx 15$ 对极，$n = \qty{200}{rpm}$。

\textbf{18.} $\mathcal P = Fu = \qty{58}{MW}$；射流功率 $\tfrac12D_mv^2 = \qty{59}{MW}$：
$98\%$；$2 \times 58 \times 0.95 = \qty{110}{MW}$。

\textbf{19.} 入口处，每千克的 $Rv = \qty{94}{m^2/s}$。出口处绝对周向速率
$u - (v - u)\cos15^\circ = \qty{1.1}{m/s}$，即 \qty{1.6}{m^2/s}。$\Gamma = D_m(94 -
1.6) = \qty{2.8}{MN.m}$——与第 16 题相同。

\textbf{20.} 当 $u \to v$ 时，相对速率、作用力与力矩
均趋于零：没有任何东西制动水轮，水轮因而超速；导流器
在一秒钟内将射流移离斗叶，同时针阀
关闭得足够慢，以避免水锤。

\textbf{21.} $v = \qty{5.7}{m/s}$，$v^2/2g = \qty{1.6}{m}$；损失 $0.015 \times 133 \times 1.6 =
\qty{3.3}{m}$；泵扬程 $200 + 3.3 + 1.6 = \qty{205}{m}$。

\textbf{22.} $\rho gD_VH = 9810 \times 40 \times 205 = \qty{80}{MW}$；电功率 \qty{91}{MW}。

\textbf{23.} $P = P_0 + \rho g(200 + 3.3) + \tfrac12\rho v^2 \approx \qty{21}{bar}$（绝对压强），
与涡轮机工况下的 \qty{20.3}{bar} 相比：此时损失
不再减去，而是加到静水头上。

\textbf{24.} $V = 40 \times 28800 = \qty{1.15e6}{m^3}$，$E = \rho gV \times 200 = \qty{2.3e12}{J}
= \qty{630}{MWh}$；耗电 $91 \times 8 = \qty{730}{MWh}$；回收 $630 \times 0.9 \times
0.95 = \qty{540}{MWh}$：往返效率 $74\%$。

\textbf{25.} 动量：喷气发动机与火箭的推力，以及
佩尔顿力。角动量：涡轮机的力矩（欧拉）。能量：
泵的扬程与功率、储能效率。质量：每一流量，
以及火箭的质量损失。
\end{solution}
